\section{Recuperation at the Discourse}
\label{sec:7.4}
There is a third instance and it is the one this book is inside. Alignment and AI ethics are, among other things, a channel through which objections to what is being built are solicited, funded, published, and answered. That channel exists, it is used, and the question section~\ref{sec:2.1.4} says to ask of any such channel is whether it has ever moved something the party collecting it did not want moved.

The argument that it has not is older than the current wave and it was made from inside. Rodrigo Ochigame, writing after working at an institution central to the field's formation, argued that the discourse of “ethical AI” was constructed with the strategic function of forestalling legally enforceable restriction — that the industry's investment in academic ethics was an investment in a substitute for regulation, and that the substitute worked by converting political questions into technical ones a company could then be trusted to answer \autocite{ochigame2019invention}. His worked example is the Partnership on AI, which this book cites repeatedly as its standing model of voluntary cross-sector coordination and treats at length in section~\ref{sec:9.1.2}: he describes its recommendations arriving reliably at conclusions compatible with the corporate agenda, with the political question of whether a technology should exist reduced to a technical question of how to build it responsibly.

\runin{The structure, tested against its own definition}
Section~\ref{sec:2.1.4}'s four features, applied here, do not give a clean verdict, so what follows is the working and not a result.

The constraint survives while its reason dies: partially. The field's founding objections were about deployment — who is scored, who is watched, who is denied — and a great deal of current work is about model properties instead, which is a shift from questions with defendants to questions with benchmarks.

Refusal is converted into a managed appearance: this is the strongest fit. An objection raised inside a lab becomes a paper, a red-team finding, a policy position, or a resignation. Each of those is a real thing that happened and none of them stopped a deployment, and the parties collecting them can point to the collection as evidence of seriousness.

Asking why becomes the deviant act: at the level of individual careers, contested. Section~\ref{sec:9.1.2}'s episode is one data point and it points the right way; there are counterexamples of people who objected loudly and prospered.

The exception cannot be filed: fits. There is no route by which a researcher's conclusion that a system should not ship becomes a decision that it does not ship, because the deciding party is the funding party and no external body holds standing.

Three of four is not a verdict of fascism at the molecular scale and I am not delivering one. It is enough to say that the mechanism has the shape, which is exactly what this book says a system should be built to notice.

\runin{What this does to the book}
The note on method that opens this book discloses that this text was written with a model made by a company the book criticizes. This section is the argument that disclosure was preparing the ground for. A book about recuperation, produced through the discourse it is describing, cannot claim to have stood outside it, and I have no procedure that would let me check whether the objections in these pages are the ones a channel like this reliably produces — the ones that are publishable, that make the author look serious, and that cost nobody anything.

What I can do is name the test and let a reader apply it. The test in section~\ref{sec:7.3} is whether disagreement has ever changed something the collecting party did not want changed. Applied here, it asks which of this book's arguments would be inconvenient to the industry that made its production possible. Chapter~\ref{sec:3} proposes systems that can refuse their owners; section~\ref{sec:6.4.1} names a vendor's participation in a targeting pipeline; section~\ref{sec:8.3.3} proposes fiscal commitments no company wants; sections~\ref{sec:9.3.3} and \ref{sec:11.2} propose binding external review of the thing chapter~\ref{sec:3} asks to be built. Those are the entries. Whether they are the real thing or the acceptable version of it is not something the person who wrote them is positioned to judge.

The one structural remedy this section can name is the one the rest of the chapter names. What the field lacks is a refusal priced by the refuser and a body with standing to receive it: an external authority that can stop a deployment, to which an objection can be filed by someone who does not depend on the objection being welcome. Section~\ref{sec:9.3.3} works out what making review binding requires, and section~\ref{sec:12.1.2} counts what a decade of the voluntary version produced. Until something of that kind exists, everything published under the heading of alignment — including this book — is being collected by a channel whose owners decide what it moves.
